\chapter{Examples and Exercises in Weak Convergence}
\label{sec:p15-weak-convergence-exercises}

This chapter completes the remaining exercise presentations in
Durrett's Section 3.2: changing densities, extreme values,
comparison of probability metrics, lattice laws, high-dimensional
sphere coordinates, and moment concentration
\cite[Exercises~3.2.1--3.2.17, pp.~123--125]{DurrettPTEAuthor2019}.
The already supplied exercises are linked at the end, without
duplicating their complete proofs.

\begin{hbdefinition}{def:p15-wce-conventions}{Laws, maxima, and normal tails}
All random variables are finite and real unless a vector domain
is stated. Weak convergence means convergence of CDFs at every
continuity point of the limiting CDF.
For independent variables \(X_1,\ldots,X_n\) with common CDF \(F\),
write \(M_n=\max_{1\leq j\leq n}X_j\), \(n\geq1\).
Set
\[
 \phi(x)=(2\pi)^{-1/2}e^{-x^2/2},\qquad
 \Phi(x)=\int_{-\infty}^x\phi(t)\,dt,\qquad
 \overline\Phi(x)=1-\Phi(x).
 \]
All logarithms are natural. The notation \(a_n\sim b_n\)
means \(a_n/b_n\to1\), with nonzero denominators eventually.
For probability measures on a common measurable space, use
\(\|\mu-\nu\|_{\mathrm{TV}}=\sup_A|\mu(A)-\nu(A)|\).
\index{weak convergence!exercise constructions}
\index{extreme-value distribution}
\end{hbdefinition}

\section{Densities and extreme values}
\label{sec:p15-wce-density-extremes}

\begin{hbtheorem}{thm:p15-wce-dyadic-densities}{proved=annex:p15.wce-dyadic-densities}{Weak convergence with failure of density convergence at every point}
For integers \(n\geq1\), let
\[
 E_n=\bigcup_{j=0}^{2^n-1}
       [j2^{-n},(j+\tfrac12)2^{-n}),\qquad
 f_n(x)=2\mathbf1_{E_n}(x).
 \]
These are probability densities, zero outside \([0,1)\),
with CDFs \(F_n\). If \(F\) is the uniform CDF on \((0,1)\),
then
\[
 \sup_x|F_n(x)-F(x)|=2^{-n-1}\longrightarrow0.
 \]
Nevertheless \(|f_n(x)-1|=1\) for every \(x\in[0,1]\)
and every \(n\). Thus no point in that whole closed interval
has density convergence to one, including its endpoints.
The laws have constant total-variation distance \(1/2\) from
the uniform law, and the densities have \(L^1\) distance one
from its density. This supplies Exercise 3.2.1.
\end{hbtheorem}

\begin{hbtheorem}{thm:p15-wce-extreme-maxima}{proved=annex:p15.wce-extreme-maxima}{All three source extreme-value limits}
For independent observations with common CDF \(F\),
\(\mathbb P(M_n\leq x)=F(x)^n\). The following full CDF limits
hold.
\begin{enumerate}
\item For \(\alpha>0\), let \(F(x)=0\) for \(x<1\) and
 \(F(x)=1-x^{-\alpha}\) for \(x\geq1\). Then
 \(M_n/n^{1/\alpha}\Rightarrow Z_\alpha\), whose CDF is
 \[
  H_\alpha(y)=
  \begin{cases}0,&y\leq0,\\ e^{-y^{-\alpha}},&y>0.\end{cases}
 \]
 This is the Fr\'echet form.
\item For \(\beta>0\), let \(F(x)=0\) for \(x<-1\),
 \(F(x)=1-|x|^\beta\) for \(-1\leq x\leq0\), and
 \(F(x)=1\) for \(x>0\). Then \(n^{1/\beta}M_n\Rightarrow W_\beta\),
 whose CDF is
 \[
  H_\beta^-(y)=
  \begin{cases}e^{-(-y)^\beta},&y<0,\\1,&y\geq0.\end{cases}
 \]
 This is the reversed Weibull form.
\item For \(F(x)=0\) at \(x<0\) and \(F(x)=1-e^{-x}\)
 at \(x\geq0\),
 \[
  M_n-\log n\Rightarrow G,\qquad
  \mathbb P(G\leq y)=e^{-e^{-y}}\quad(y\in\mathbb R).
 \]
 This is the Gumbel, or extreme-value double-exponential, law.
\end{enumerate}
All limiting CDFs are continuous and normalized.
For a probability level \(0<u<1\), the exact respective
unscaled maximum quantiles are
\[
 (1-u^{1/n})^{-1/\alpha},\qquad
 -(1-u^{1/n})^{1/\beta},\qquad
 -\log(1-u^{1/n}).
 \]
Their stated scalings converge to the corresponding limiting
quantiles. These are all three cases of Exercise 3.2.2.
\end{hbtheorem}

Durrett closes Exercise~3.2.2 with a naming and a historical
remark \cite[p.~123]{DurrettPTEAuthor2019}. He calls the three
limits above the extreme value distributions, and the third the
double exponential or Gumbel distribution. He then records that
Gnedenko (1943) found necessary and sufficient conditions on \(F\)
for affinely normalized maxima \((M_n-b_n)/a_n\), with deterministic
\(a_n>0\) and \(b_n\), to converge to these limits, and he refers to
Resnick (1987) for a recent treatment. The source states this as an
attribution only: it gives neither the conditions nor their proof.
The three models of \cref{thm:p15-wce-extreme-maxima} are particular
cases with \(b_n=0\), \(a_n=n^{1/\alpha}\); \(b_n=0\),
\(a_n=n^{-1/\beta}\); and \(b_n=\log n\), \(a_n=1\).

\begin{hbtheorem}{thm:p15-wce-normal-maxima}{proved=annex:p15.wce-normal-maxima}{The complete Gaussian maximum normalization}
For every fixed \(\theta\in\mathbb R\),
\[
 \frac{\overline\Phi(x+\theta/x)}{\overline\Phi(x)}
                                    \longrightarrow e^{-\theta}
                                      \quad(x\to\infty).
 \]
For \(n\geq3\), let \(b_n>0\) be the unique solution of
\(\overline\Phi(b_n)=1/n\).
For i.i.d. standard normal \(X_j\),
\[
 \mathbb P\{b_n(M_n-b_n)\leq y\}\longrightarrow e^{-e^{-y}}
                                               \quad(y\in\mathbb R).
 \]
The centering sequence satisfies
\[
 b_n\sim\sqrt{2\log n},\qquad
 b_n^2=2\log n-\log\log n-\log(4\pi)+o(1),
 \qquad
 \frac{M_n}{\sqrt{2\log n}}\longrightarrow1
                                     \quad\text{in probability}.
 \]
The tail equation has no finite solution at \(n=1\) and gives
\(b_2=0\); restricting its use to \(n\geq3\) removes these
irrelevant initial exceptions. This supplies all parts of Exercise 3.2.3.
\end{hbtheorem}

\section{Metrics and integer-valued laws}
\label{sec:p15-wce-metrics-lattices}

\begin{hbtheorem}{thm:p15-wce-bounded-distance}{proved=annex:p15.wce-bounded-distance}{The source comparison with the Ky Fan metric}
For variables on one probability space, let
\[
 \alpha(X,Y)=\inf\{\varepsilon\geq0:
               \mathbb P(|X-Y|>\varepsilon)\leq\varepsilon\},
 \qquad
 \beta(X,Y)=\mathbb E\frac{|X-Y|}{1+|X-Y|}.
 \]
The already proved Ky Fan metric has attained infimum and
\(0\leq\alpha\leq1\).
The function \(\beta\) is a metric on variables modulo
almost-sure equality, requiring no moments.
If \(a=\alpha(X,Y)\), then
\[
 \frac{a^2}{1+a}\leq\beta(X,Y)
       \leq a+(1-a)\frac a{1+a}=\frac{2a}{1+a}.
 \]
These bounds imply
\(\beta(X_n,X)\to0\) exactly when \(X_n\to X\) in probability.
The lower bound can be attained for every \(0\leq a\leq1\);
the upper bound is approached by explicit finite-valued examples.
This proves Exercise 3.2.8 and its needed Exercise 2.3.6 premise.
\index{Ky Fan metric!bounded-distance comparison}
\end{hbtheorem}

\begin{hbtheorem}{thm:p15-wce-integer-laws}{proved=annex:p15.wce-integer-laws}{Weak convergence of integer laws is convergence of all their masses}
Let \(X_n\), \(1\leq n\leq\infty\), be integer-valued.
Then
\[
 X_n\Rightarrow X_\infty
 \quad\Longleftrightarrow\quad
 \mathbb P(X_n=m)\longrightarrow\mathbb P(X_\infty=m)
                                      \quad(m\in\mathbb Z).
 \]
In either case the mass functions converge in \(\ell^1(\mathbb Z)\)
and the laws converge in total variation.
The limiting mass function is the law of the specified
integer-valued \(X_\infty\), so its total mass is one.
This normalization is essential to Exercise 3.2.11.
\end{hbtheorem}

\section{Uniform surface measure and Gaussian coordinates}
\label{sec:p15-wce-sphere}

\begin{hbdefinition}{def:p15-wce-spherical-measure}{Surface area defined independently of Gaussian coordinates}
For integer \(d\geq1\), let \(\lambda_d\) be product Lebesgue
measure on \(\mathbb R^d\), let
\(\|x\|=(\sum_{j=1}^d x_j^2)^{1/2}\), and put
\[
 B_d=\{x:\|x\|\leq1\},\qquad
 \mathbb S^{d-1}=\{x:\|x\|=1\},\qquad v_d=\lambda_d(B_d).
 \]
An orthogonal map is a real linear map \(O\) with
\(O^{\mathsf T}O=I\); it preserves this norm.
For a Borel set \(A\subseteq\mathbb S^{d-1}\), define its cone
\[
 C(A)=\{ru:u\in A,\ 0<r\leq1\},\qquad
 \sigma_d(A)=\lambda_d(C(A))/v_d.
 \]
For radius \(r>0\), spherical surface area on the patch \(rA\)
is defined by the radial-shell limit
\[
 \operatorname{area}_r(rA)=
 \lim_{\varepsilon\downarrow0}
 \frac{\lambda_d(\{tu:u\in A,\ r<t\leq r+\varepsilon\})}
      {\varepsilon}.
 \]
The next result proves existence of this limit and its normalization.
Uniform surface probability at radius \(r\) means this area
divided by total area. This geometric definition is made before
introducing the Gaussian representation.
At \(d=1\), the sphere consists of two points; its area gives
each point mass one before normalization.
\index{surface measure!sphere}\index{polar integration!cone construction}
\end{hbdefinition}

\begin{hbtheorem}{thm:p15-wce-spherical-area}{proved=annex:p15.wce-spherical-area}{The cone construction and polar integral}
One has \(0<v_d<\infty\), and \(\sigma_d\) is an
orthogonally invariant Borel probability on \(\mathbb S^{d-1}\).
For every \(r>0\),
\[
 \operatorname{area}_r(rA)=d r^{d-1}v_d\,\sigma_d(A).
 \]
In particular normalized surface probability at radius \(r\)
is the image of \(\sigma_d\) under \(u\mapsto ru\).
For every nonnegative Borel function \(h\),
\[
 \int_{\mathbb R^d}h(x)\,d\lambda_d(x)
 =d v_d\int_0^\infty\int_{\mathbb S^{d-1}}
                       h(ru)\,d\sigma_d(u)\,r^{d-1}\,dr.
 \]
The identity also holds for integrable signed functions.
Its proof uses rectangle measure uniqueness and scalar dilations,
not an assumed general multivariable change-of-variables theorem.
The constants are
\[
 v_d=\frac{\pi^{d/2}}{\Gamma(d/2+1)},\qquad
 d v_d=\frac{2\pi^{d/2}}{\Gamma(d/2)}.
 \]
Here \(\Gamma\) is the gamma integral of
Definition~\ref{def:p15-gamma-density}.
\end{hbtheorem}

\begin{hbtheorem}{thm:p15-wce-sphere-coordinate}{proved=annex:p15.wce-sphere-coordinate}{The full surface-coordinate exercise}
Let \(X^{(n)}\) have uniform surface probability on the
sphere of radius \(\sqrt n\) in \(\mathbb R^n\).
Then its first coordinate satisfies
\[
                         X^{(n)}_1\Rightarrow N(0,1).
 \]
More precisely, on a space with independent standard normals
\(Y_1,Y_2,\ldots\), set
\[
 D_n=\frac1n\sum_{j=1}^nY_j^2,\qquad
 Z^{(n)}=\frac{(Y_1,\ldots,Y_n)}{\sqrt{D_n}}
                                      \quad\text{when }D_n>0.
 \]
On \(\{D_n=0\}\), set \(Z^{(n)}=\sqrt n(1,0,\ldots,0)\).
Then \(Z^{(n)}\) has exactly the prescribed uniform surface law,
and \(Z^{(n)}_1\to Y_1\) almost surely.
This coupling proves the asserted convergence of the coordinate laws;
it does not impose that coupling on an arbitrary original sequence.
For every \(n\geq1\), the coordinate has
\[
 \mathbb EX^{(n)}_1=0,\qquad
 \mathbb E(X^{(n)}_1)^2=1,\qquad
 \mathbb E(X^{(n)}_1)^4=\frac{3n}{n+2}.
 \]
This supplies the geometric and probabilistic premises of Exercise 3.2.15.
\end{hbtheorem}

\section{Moment concentration and uniform lower bounds}
\label{sec:p15-wce-moments}

\begin{hbtheorem}{thm:p15-wce-moment-concentration}{proved=annex:p15.wce-moment-concentration}{Two distinct positive moments force a constant limit}
Let \(Y_n\geq0\) and \(0<\alpha<\beta\).
If
\[
                       \mathbb EY_n^\alpha\to1,\qquad
                       \mathbb EY_n^\beta\to1,
 \]
then \(Y_n\to1\) in probability. No independence is required.
For \(p=\beta/\alpha>1\), define
\(D(z)=z^p-pz+p-1\), \(z\geq0\).
For \(0<\varepsilon<1\), a quantitative bound for all
sufficiently large \(n\) is
\[
 \mathbb P(|Y_n-1|>\varepsilon)
 \leq
 \frac{\mathbb EY_n^\beta-p\mathbb EY_n^\alpha+p-1}
 {\min\{D((1-\varepsilon)^\alpha),D((1+\varepsilon)^\alpha)\}}.
 \]
The denominator is strictly positive, and all subtracted
expectations are finite at those sufficiently large indices.
This is Exercise 3.2.16 for its whole real exponent range.
\end{hbtheorem}

\begin{hbtheorem}{thm:p15-wce-fourth-moment}{proved=annex:p15.wce-fourth-moment}{The source uniform probability bound}
For \(K\geq1\), if \(\mathbb EX^2=1\) and
\(\mathbb EX^4\leq K<\infty\), then for \(0\leq y<1\),
\[
                 \mathbb P(|X|>y)\geq\frac{(1-y^2)^2}{K}>0.
 \]
For \(y<0\) the probability is one.
If \(K<1\), no variable can meet the two moment assumptions,
so any positive lower-bound constant, for example one, works vacuously.
At \(y=0\), the bound \(1/K\) can be attained.
This supplies every parameter case of Exercise 3.2.17.
\end{hbtheorem}

\begin{hbtheorem}{thm:p15-wce-boundaries}{proved=annex:p15.wce-boundaries}{Why the exercise hypotheses matter}
The following explicit examples separate the assumptions.
\begin{enumerate}
\item Giving all Pareto observations the same variable preserves their
 marginal CDF but makes their normalized maximum tend to zero,
 instead of the independent Fr\'echet limit.
\item \(X_n=1/n\) converges weakly to zero while its point
 mass at zero stays zero. The lattice hypothesis is needed
 for the stated mass-convergence necessity.
 Integer constants \(X_n=n\) have point masses tending to
 zero at every integer, but that zero sequence is not a probability law.
\item One positive moment tending to one is insufficient:
 a fixed variable taking \(0\) and \(2^{1/\alpha}\) with equal
 probabilities has \(\mathbb EY^\alpha=1\) and does not
 concentrate at one. Nonnegativity is also substantive:
 \(Y_n=-1\), with exponents \(2,4\), has both moments one
 and does not converge to \(+1\).
\item Without a common fourth-moment bound, variables taking
 \(\sqrt n\) with probability \(1/n\) and zero otherwise
 have second moment one while
 \(\mathbb P(|X_n|>y)\to0\) for every fixed \(0\leq y<1\).
 At threshold \(y=1\), the fair-sign law has second and
 fourth moments one but zero strict-exceedance probability.
\end{enumerate}
\end{hbtheorem}

\section{The already supplied source exercises}
\label{sec:p15-wce-existing-routes}

Exercises 3.2.4, 3.2.5 and 3.2.10 are supplied by
Theorem~\ref{thm:p15-criteria-weak-integration}: weak Fatou,
the unbounded integration criterion, and the finite quantile
approximation with its exact uniform CDF error.
Exercises 3.2.6, 3.2.7 and 3.2.12 are supplied by
Theorem~\ref{thm:p15-weak-metric-premises}, including attained
Ky Fan distance and both probability/constant-limit implications.
Exercise 3.2.9 is
Theorem~\ref{thm:p15-continuous-cdf-uniform}.
For Exercises 3.2.13 and 3.2.14, combine the constant-limit
bridge just cited with Theorem~\ref{thm:p15-clt-perturbation};
the latter supplies sums, products and the source's unrestricted
real constant version without an independence assumption.
For 3.2.13, the bridge gives \(Y_n-c\to0\) in probability;
use that residual in the sum lemma. Its stated consequence
uses the residual \(Z_n-X_n\).
For 3.2.14, the bridge gives \(Y_n\to c\) in probability;
set \(A_n=Y_n\) and \(B_n=0\) in the product assertion.

\begin{proofboundary}{def:p15-wce-completion}{The completed exercise list}
Together these routes supply all seventeen Section 3.2 exercises,
including the closing remark of Exercise~3.2.2, which is stated
after \cref{thm:p15-wce-extreme-maxima} in the form the source gives
it: the naming of the extreme value and Gumbel distributions and the
attribution of the necessary and sufficient conditions to Gnedenko
(1943), with Resnick (1987) as the recent treatment. The source
attributes those conditions without stating or proving them, and the
three explicit model limits above are not a proof of them.
\end{proofboundary}
